\documentclass[11pt]{article}
% ======================================================================
%  Preambule commun - Sujets Bac Serie A (Mathematiques)
%  Style sobre : noir et blanc, sans couleurs, sans filigrane,
%  sans en-tetes ni pieds de page decoratifs.
% ======================================================================
\usepackage[utf8]{inputenc}
\usepackage[T1]{fontenc}
\usepackage{lmodern}
\usepackage{textcomp}
\usepackage{amsmath,amssymb}
\usepackage{enumitem}
\usepackage{array}
\usepackage[a4paper,margin=2.3cm]{geometry}
\usepackage{fancyhdr}
\usepackage{tikz}

% Symbole degre robuste + justification souple
\DeclareUnicodeCharacter{00B0}{\ensuremath{^\circ}}
\tolerance=2000
\emergencystretch=2em
\frenchspacing

% En-tete / pied de page minimalistes (numero de page seul)
\pagestyle{fancy}
\fancyhf{}
\renewcommand{\headrulewidth}{0pt}
\fancyfoot[C]{\thepage}

% Macros mathematiques
\newcommand{\R}{\mathbb{R}}
\newcommand{\N}{\mathbb{N}}
\newcommand{\Z}{\mathbb{Z}}
\newcommand{\vi}{\vec{\imath}}
\newcommand{\vj}{\vec{\jmath}}

% Titres de blocs
\newcommand{\exo}[2]{\bigskip\noindent{\large\bfseries Exercice #1}\hfill\textit{#2}\par\nopagebreak\smallskip}
\newcommand{\partie}[1]{\medskip\noindent{\bfseries Partie #1}\par\nopagebreak\smallskip}
\newcommand{\entetesujet}[1]{\begin{center}{\Large\bfseries #1}\end{center}\vspace{0.3em}\hrule\vspace{1.1em}}

% Questions numerotees : niveau 1 = chiffres gras, niveau 2 = a. b. c.
\newlist{questions}{enumerate}{1}
\setlist[questions]{label=\textbf{\arabic*.},leftmargin=1.8em,itemsep=3pt,topsep=3pt}
\newlist{sousq}{enumerate}{1}
\setlist[sousq]{label=\textbf{\alph*.},leftmargin=1.8em,itemsep=2pt,topsep=2pt}


\begin{document}

\entetesujet{Sujet bac 2019 -- Série A}

\exo{1}{8 points}
\partie{1}
\begin{questions}
  \item \begin{sousq}
    \item Décomposer les nombres $1\,155$ et $300$ en produit de puissances de nombres premiers.
    \item En déduire le $\mathrm{PGCD}(1155\,;300)$.
  \end{sousq}
  \item Rendre irréductible la fraction $\dfrac{1\,155}{300}$.
  \item La fraction $\dfrac{1\,155}{300}$ est-elle décimale ?

  \underline{NB.} On utilisera la propriété suivante : une fraction est décimale si le dénominateur de sa fraction irréductible ne contient que les puissances de 2 et/ou de 5.
\end{questions}

\partie{2}
\begin{questions}
  \item On donne le polynôme $p(x) = -2x^3 + 9x^2 - 7x - 6$. En considérant $x_0 = 2$, une solution évidente de l'équation $p(x) = 0$, déterminer les réels $a$ et $b$ tels que $p(x) = (x - 2)(ax^2 + 5x + b)$.
  \item Considérons $a = -2$ et $b = 3$.
  \begin{sousq}
    \item Résoudre dans $\R$ l'équation $-2x^2 + 5x + 3 = 0$.
    \item En déduire les solutions de l'équation $p(x) = 0$ dans $\R$.
  \end{sousq}
  \item Résoudre dans $\R$ l'équation $(E) : -(e^{t})^3 + 9(e^{t})^2 - 7\,e^{t} - 6 = 0$.
  \item \begin{sousq}
    \item Trouver la valeur de $Y$ dans le système $S$ pour $X = -1$. On donne $S : \begin{cases} -2X + Y = 4 \\ 3X + 2Y = 1 \end{cases}$
    \item En déduire, dans $\R^2$, les solutions du système $(S') : \begin{cases} -2\ln x + e^{y} = 4 \\ 3\ln x + 2\,e^{y} = 1 \end{cases}$
  \end{sousq}
\end{questions}

\exo{2}{8 points}
On considère la fonction numérique $f$ de la variable réelle $x$ définie par : $f(x) = \dfrac{ax + 3}{x + 1}$ avec $a$ un réel. On désigne par $(\mathcal{C}_f)$ la courbe représentative de $f$ dans un plan $(\mathcal{P})$ muni d'un repère orthonormé $(O,\vi,\vj)$. Unité graphique : 1 cm.
\begin{questions}
  \item Déterminer l'ensemble de définition $E_f$ de $f$ sous forme d'une réunion d'intervalles.
  \item Trouver le réel $a$ pour que la courbe de $f$ passe par $A(1\,;1)$.
  \item Dans la suite, on posera $a = -1$.
  \begin{sousq}
    \item Calculer la limite de $f$ en $-\infty$ et la limite de $f$ en $-1$ à gauche.
    \item Donner l'interprétation graphique des limites à l'infini et en $-1$ calculées au a.
  \end{sousq}
  \item \begin{sousq}
    \item Calculer la fonction dérivée $f'$ de $f$.
    \item Préciser le signe de $f'(x)$ pour $x \in E_f$.
    \item Préciser le sens de variation de $f$ sur chacun des intervalles de $E_f$.
    \item Dresser le tableau de variation de $f$. On donne $\lim\limits_{x\to+\infty} f(x) = -1$ et $\lim\limits_{x\to-1^+} f(x) = +\infty$.
  \end{sousq}
  \item Construire, dans le plan $(\mathcal{P})$ muni du repère $(O,\vi,\vj)$, la courbe $(\mathcal{C})$ et les droites $(\mathcal{D}_1) : x = -1$ et $(\mathcal{D}_2) : y = -1$ qui sont les asymptotes à la courbe.
\end{questions}

\exo{3}{4 points}
Chaque année le taux de natalité est calculé en divisant le nombre de naissances par le nombre d'habitants d'un pays. Dans un pays en développement, on a constaté lors d'une étude statistique que le taux de natalité était en baisse. Voici ci-dessous le tableau des résultats constatés.
\begin{center}\renewcommand{\arraystretch}{1.4}
\begin{tabular}{|c|c|c|c|c|c|}
\hline
Rang de l'année ($x_i$) & 4 & 6 & 8 & 10 & 12 \\
\hline
Taux de natalité en \% ($y_i$) & 1{,}5 & 3{,}1 & 2{,}8 & 2{,}5 & 1{,}3 \\
\hline
\end{tabular}
\end{center}
\begin{questions}
  \item Construire le nuage des points $M_i(x_i, y_i)$ associé à cette série statistique dans un plan muni d'un repère orthogonal. On prendra : sur l'axe $(x'Ox)$, 1 cm pour une année ; sur l'axe $(y'Oy)$, 1 cm pour 1 \%.
  \item On désigne par $G_1(x_{G_1}, y_{G_1})$ le point moyen du sous-nuage associé au tableau suivant :
  \begin{center}\renewcommand{\arraystretch}{1.4}
  \begin{tabular}{|c|c|c|}\hline Rang de l'année ($x_i$) & 4 & 6 \\\hline Taux de natalité en \% ($y_i$) & 1{,}5 & 3{,}1 \\\hline\end{tabular}
  \end{center}
  De même $G_2(x_{G_2}, y_{G_2})$ le point moyen du sous-nuage associé au tableau ci-après :
  \begin{center}\renewcommand{\arraystretch}{1.4}
  \begin{tabular}{|c|c|c|c|}\hline Rang de l'année ($x_i$) & 8 & 10 & 12 \\\hline Taux de natalité en \% ($y_i$) & 2{,}8 & 2{,}5 & 1{,}3 \\\hline\end{tabular}
  \end{center}
  \begin{sousq}
    \item Calculer les coordonnées de $G_1$ et de $G_2$.
    \item Déterminer, sous la forme $y = ax + b$ (avec $a$ et $b$ des réels), l'équation de la droite $(G_1G_2)$.
    \item On suppose que le taux de natalité $y$ est donné en fonction du rang de l'année par la formule : $y = -\dfrac{1}{50}x + \dfrac{12}{5}$. Déterminer le taux de natalité à la quatorzième ($14^{\text{e}}$) année.
  \end{sousq}
\end{questions}

\end{document}
